\chapter{Usage au quotidien}
\label{ch:usage}

Ce chapitre parcourt les flux que vous lancerez des dizaines de
fois une fois la bibliothèque opérationnelle : cataloguer un nouveau
volume, chercher dans la collection, déplacer un volume, le prêter,
le récupérer, et auditer ce qui est sur les étagères.

\section{La page d'accueil}\label{sec:home}

\begin{figure}[ht]
  \centering
  \includegraphics[width=0.95\textwidth]{01-home-desktop.png}
  \caption{La page d'accueil sur une installation réelle~:
    champ de recherche, puces de filtres par genre, compteurs
    «~À traiter~», totaux «~Aperçu de la collection~» et bandeau
    d'ajouts récents avec les jaquettes récupérées via la chaîne
    de fournisseurs de métadonnées.}
  \label{fig:home}
\end{figure}

La page d'accueil est le tableau de bord. De haut en bas :

\begin{itemize}
  \item \textbf{Champ de scan.}\index{champ de scan} Un champ
        focalisé qui accepte un ISBN, un code V ou un code L. La
        couche scanner-guard route l'entrée vers la bonne page
        selon le préfixe.
  \item \textbf{Bandeau d'indicateurs.} Ajouts récents, prêts en
        retard, séries avec lacunes, volumes non rangés — chaque
        clic ouvre la liste filtrée correspondante.
  \item \textbf{Activité récente.} Les sept derniers jours de
        catalogage et de retours.
\end{itemize}

\section{Cataloguer un nouveau titre via code-barres}\label{sec:catalog-scan}

\begin{enumerate}
  \item Focalisez le champ de scan de la page d'accueil (il est
        sélectionné par défaut au chargement). Scannez l'ISBN /
        EAN-13 du dos de l'ouvrage.
  \item mybibli redirige vers la page de catalogage, qui affiche un
        titre squelette avec le code-barres scanné. Une tâche
        d'arrière-plan lance des requêtes contre la chaîne de
        fournisseurs de métadonnées.
  \item En quelques secondes, le squelette se remplit : titre,
        contributeurs, éditeur, année, image de couverture. Un
        petit indicateur confirme quel fournisseur a résolu le
        titre.
  \item Choisissez un \textbf{lieu}\index{lieu} dans le menu
        déroulant (ou scannez le code L du rayon de destination) et
        un \textbf{état de volume}\index{état de volume}
        (typiquement \emph{Neuf}). Cliquez \textbf{Ajouter un
        volume}.
  \item Imprimez et collez une étiquette de code V sur le volume.
\end{enumerate}

Si le pipeline de métadonnées échoue à résoudre l'ISBN, le titre
apparaît avec un badge \emph{Non résolu}. Vous pouvez modifier les
champs manuellement et le volume est créé malgré tout.

\begin{tip}
Le catalogage en série marche le mieux quand la planche de codes V,
la douchette et l'imprimante sont à portée de main. La plupart du
temps est consacré à scanner, pas à attendre mybibli — la
récupération de métadonnées est asynchrone, vous n'avez donc pas à
attendre qu'elle se termine avant de passer au volume suivant.
\end{tip}

\subsection*{Imprimer de nouvelles planches d'étiquettes}

Avant d'imprimer une nouvelle planche d'étiquettes V ou L, consultez
la petite ligne en haut de la page de catalogage (depuis la 1.12.0,
Librarian et plus) : \emph{« Étiquettes utilisées jusqu'à
V0142~·~L0037 — imprimez les nouvelles planches à partir des numéros
suivants. »} Les numéros incluent les éléments actuellement à la
corbeille — une étiquette imprimée et collée un jour sur un livre ne
sera jamais reproposée, même si ce volume a été supprimé depuis.

\section{Rechercher dans la collection}\label{sec:search}

La page \texttt{/search} accepte :

\begin{itemize}
  \item \textbf{titre} complet ou partiel ;
  \item \textbf{nom de contributeur} (matche l'auteur par défaut ;
        les autres rôles de contributeur sont scorés moins fort) ;
  \item \textbf{ISBN / EAN-13} — match exact ;
  \item \textbf{code V} ou \textbf{code L} — match exact ;
  \item \textbf{code Dewey} (pour les bibliothèques qui utilisent
        Dewey).
\end{itemize}

Les résultats sont triés par pertinence. Cliquez un résultat pour
ouvrir la page du titre ; cliquez un contributeur pour filtrer le
catalogue sur ses œuvres ; cliquez le libellé d'un lieu pour voir
chaque volume actuellement à cet endroit.

\begin{tip}
Vous pouvez « rechercher » un seul volume en scannant son code V
depuis n'importe quelle page qui dispose du champ de scan d'accueil
— y compris la page d'accueil, \texttt{/search} et \texttt{/loans}.
Plus rapide que la frappe.
\end{tip}

\subsection{Recherches sauvegardées}\label{sec:saved-searches}

Les recherches récurrentes (« BD sans couverture », « titres non
classés », « tout ce qui est en retard ») peuvent être sauvegardées
une fois puis relancées d'un seul clic. Sur la page d'accueil,
composez la recherche voulue — saisissez une requête, cliquez une
puce de filtre, choisissez un tri — puis ouvrez le menu déroulant
\textbf{Recherches sauvegardées} à côté du contrôle de tri et utilisez
\emph{Sauver la recherche courante} pour la nommer et l'enregistrer.

Les recherches sauvegardées sont \textbf{partagées par toute
l'instance} : tous les bibliothécaires voient et peuvent lancer le
même ensemble. Pour en lancer une, ouvrez le menu et cliquez son nom
— la page d'accueil se recharge avec cette requête, ce filtre et ce
tri appliqués. Chaque entrée propose aussi \emph{Renommer} et
\emph{Supprimer} (confirmation par fenêtre modale). La suppression est
réversible : sauvegarder une nouvelle recherche du même nom restaure
l'ancienne.

\begin{note}
Les recherches sauvegardées capturent les critères de navigation
existants (requête, une puce de filtre, colonne et sens de tri) — ce
n'est pas un constructeur de requêtes libre. Le contrôle n'est visible
que pour les bibliothécaires et les administrateurs.
\end{note}

\section{Étiquettes de gestion}\label{sec:management-labels}

Un genre dit ce qu'un livre \emph{est}. Un état dit comment un
exemplaire a \emph{vieilli}. Ni l'un ni l'autre ne dit « je ne suis pas
sûr que cette notice soit juste, il faudra y revenir » --- et c'est la
phrase dont un catalogue vivant a le plus besoin. Depuis la 1.18.0, les
\emph{étiquettes de gestion}\index{étiquette!de gestion} la portent.

Ce ne sont pas les étiquettes de dos de la section~\ref{sec:labels} :
rien ne s'imprime, rien ne se scanne, et un visiteur anonyme ne les voit
jamais. Elles existent pour la personne qui tient la bibliothèque.

\subsection*{Définir le vocabulaire}

Un administrateur crée le vocabulaire une fois, dans
\emph{Administration \textrightarrow{} Données de référence
\textrightarrow{} Étiquettes} : « À vérifier », « À relire », « Reliure
abîmée », « Prêté sans fiche » --- ce que votre pratique réclame. On les
renomme et on les supprime au même endroit, et comme toute autre valeur
de référence elles ne sont jamais traduites : ce que vous saisissez est
ce que tout le monde voit, quelle que soit la langue de l'interface.

Gardez la liste courte. Un vocabulaire de six étiquettes s'utilise ;
un vocabulaire de trente est une seconde taxonomie de genres que
personne n'entretient.

\subsection*{Poser une étiquette}

Un bibliothécaire pose les étiquettes depuis la section
\emph{Étiquettes} d'une page de titre ou d'une page de volume, en
choisissant dans ce vocabulaire commun. La distinction compte :

\begin{itemize}
  \item sur un \textbf{titre}, l'étiquette parle de la notice --- « les
        métadonnées ont l'air fausses » ;
  \item sur un \textbf{volume}, elle parle de cet exemplaire physique
        --- « celui-ci a la reliure abîmée », ce qui ne dit rien des
        deux autres exemplaires du rayon.
\end{itemize}

\subsection*{Retrouver ce qui porte une étiquette}

\texttt{/labels} liste le vocabulaire avec un compte en regard de chaque
entrée. En ouvrir une montre ce qui la porte, en deux sections --- les
titres au-dessus des volumes --- chaque ligne menant à sa page de
détail. C'est tout l'intérêt : une étiquette d'où l'on ne peut pas
revenir est un mot dans un tiroir.

\begin{note}
Les étiquettes sont internes. Les gestionnaires n'exécutent même pas les
requêtes pour un visiteur anonyme : rien ne fuit par le catalogue
public --- pas même le fait qu'une étiquette existe.
\end{note}

\section{Déplacer un volume vers un autre lieu}\label{sec:move-volume}

\begin{enumerate}
  \item Ouvrez la page du volume (recherche par code V ou clic
        depuis la page du titre).
  \item Cliquez \textbf{Déplacer}.
  \item Scannez le code L du lieu de destination, ou choisissez-le
        dans le menu déroulant.
  \item Confirmez.
\end{enumerate}

mybibli enregistre le déplacement dans l'historique de lieu du
volume, vous permettant de retracer où un livre a vécu au fil du
temps.

\section{Prêts}\label{sec:loans}

\subsection*{Prêter un volume}

\begin{enumerate}
  \item Ouvrez \texttt{/loans} et cliquez \textbf{Nouveau prêt}.
  \item Scannez ou choisissez l'\textbf{emprunteur}\index{emprunteur}.
        Si l'emprunteur n'existe pas encore, cliquez \emph{Créer
        l'emprunteur} en ligne.
  \item Scannez ou choisissez le \textbf{volume} (code V).
  \item Réglez optionnellement une \textbf{date de retour
        prévue} ; sinon mybibli la calcule à partir du seuil de
        retard (chapitre~\ref{ch:configuration}).
  \item Confirmez.
\end{enumerate}

Le lieu actuel du volume est inscrit sur la ligne de prêt pour
qu'un retour ultérieur puisse le restaurer au bon rayon.

\subsection*{Retourner un volume}

\begin{enumerate}
  \item Ouvrez \texttt{/loans} et trouvez le prêt actif (filtre
        \emph{Actifs} ou scan du code V).
  \item Cliquez \textbf{Retourner}.
  \item Choisissez l'état post-retour (\emph{Bon}, \emph{Usé},
        \dots).
  \item Confirmez. Le volume retourne automatiquement à son lieu
        d'avant prêt.
\end{enumerate}

\subsection*{Prêts en retard}

Les prêts qui dépassent le seuil de retard passent au rouge sur la
page \texttt{/loans} et sur l'indicateur \emph{En retard} de la page
d'accueil. \texttt{/loans?filter=overdue} donne une liste focalisée.

\section{Audit des rayons}\label{sec:audit}

Un \emph{audit}\index{audit} c'est le rituel « je parcours les
étagères et je confirme que tout est bien à la place que la base
indique ». Depuis la 1.7.x, mybibli porte ce flux au lieu de vous
laisser l'improviser autour de la page de recherche.

L'unité de travail est un drapeau posé sur un volume :
\emph{à contrôler}. On lève le drapeau, on parcourt le rayon, et on
l'abaisse à mesure que chaque exemplaire est confirmé.

\subsection*{Lever le drapeau}

\begin{itemize}
  \item \textbf{Un volume} — depuis sa page de détail,
        \emph{Marquer à contrôler}. Utile quand un lecteur signale que
        quelque chose cloche, ou qu'un scan a fait apparaître un volume
        là où il n'aurait pas dû être.
  \item \textbf{Un rayon entier} — depuis la page d'un lieu de
        rangement, \emph{Marquer tous les volumes à contrôler}. C'est
        ainsi qu'une passe planifiée commence : on marque le rayon le
        dimanche soir, on le parcourt dans la semaine.
\end{itemize}

Un volume marqué porte une pastille \emph{À contrôler} partout où il
est listé : le drapeau se voit sans ouvrir la vue d'audit.

\subsection*{Parcourir le rayon}

Le tableau de bord d'accueil affiche un indicateur \emph{Volumes à
contrôler} avec son compte et un lien ; \texttt{/audit} est la liste
elle-même, triée par lieu puis par code V — l'ordre dans lequel on
marche. Chaque ligne porte un bouton \emph{Contrôlé} : on l'actionne
quand on tient l'exemplaire, et la ligne quitte la liste.

Deux situations se dénouent en chemin :

\begin{itemize}
  \item Un volume sur le rayon que la liste n'attend pas — scannez-le ;
        mybibli signale qu'il appartient ailleurs et vous le déplacez
        (section~\ref{sec:move-volume}).
  \item Un volume que la liste attend et que le rayon ne porte pas —
        laissez le drapeau levé. Une liste d'audit inachevée est une
        liste de choses à faire, pas un échec : c'est le seul relevé de
        ce que vous n'avez pas retrouvé.
\end{itemize}

\emph{Effacer tous les drapeaux} vide la liste d'un coup. À utiliser
quand on abandonne une passe, pas pour en terminer une — un drapeau
effacé sans avoir le volume en main enregistre une vérification qui
n'a pas eu lieu.

\begin{tip}
Un audit trimestriel attrape les volumes mal rangés tôt — plus un
livre traîne au mauvais rayon, plus il est difficile de se rappeler
où il aurait dû être.
\end{tip}

\section{Rangement après retours}\label{sec:shelving}

Après une fournée de retours, les volumes s'empilent près du poste
de scan. Le flux de \emph{rangement}\index{rangement} de mybibli
accélère la remise en rayon :

\begin{enumerate}
  \item Prenez un volume et scannez son code V.
  \item mybibli affiche la page du volume avec le lieu d'origine
        bien en évidence.
  \item Portez le volume au rayon indiqué. Scannez le code L du
        rayon pour confirmer.
  \item Recommencez.
\end{enumerate}

La confirmation par code L est optionnelle mais recommandée — elle
laisse une entrée fraîche dans l'historique du volume sur laquelle
le prochain audit pourra s'appuyer.

\subsection{Annuler une action de scan}\label{sec:scan-undo}

Le travail au code-barres va vite, et une erreur de scan arrive
parfois — un volume rangé au mauvais lieu, ou un lieu de lot activé
par mégarde. Lorsque vous rangez un volume ou activez un lieu de lot
sur la page de catalogage, le message de confirmation affiche un
bouton \textbf{Annuler}. Cliquez-le dans les \textbf{30 secondes} et
mybibli annule cette dernière action :

\begin{itemize}
  \item Annuler un rangement remet le volume à son lieu précédent
        (ou le laisse non rangé s'il n'avait pas de lieu auparavant).
  \item Annuler l'activation d'un lieu de lot rétablit le lieu actif
        précédent.
\end{itemize}

Seule la \emph{dernière} action peut être annulée, et le bouton
disparaît une fois la fenêtre de 30 secondes écoulée. L'annulation ne
touche que le changement de lieu — elle ne supprime jamais un titre
ou un volume que vous venez de créer. Passé le délai, déplacez le
volume de la manière habituelle (\autoref{sec:move-volume}).

\section{Wish list}\label{sec:wishlist}

La \emph{wish list}\index{wish list} est une liste séparée des
livres que vous aimeriez acheter, distincte du catalogue possédé.
Elle vit sous \texttt{/wishlist}\index{wishlist@/wishlist} (voir
la barre de navigation) et est disponible pour les Bibliothécaires
et les Admins.

\subsection{Ajouter une entrée}

Le formulaire \emph{Ajouter à la wish list} propose deux modes~:

\begin{itemize}
  \item \textbf{Par ISBN.} Scannez ou saisissez un ISBN à 13
        chiffres. mybibli interroge la même chaîne de fournisseurs
        de métadonnées que le catalogue, puis affiche une carte de
        prévisualisation (titre, auteur, éditeur, année,
        couverture). Cliquez sur \emph{Ajouter à la wish list}
        pour enregistrer. Les images de couverture pointent
        directement sur l'URL du fournisseur résolu (Open Library
        Covers ou Google Books par exemple)~; le téléchargement
        local de la couverture est sur la roadmap.
  \item \textbf{Saisie libre.} À utiliser quand vous n'avez pas
        d'ISBN sous la main (titre auto-édité, recommandation
        notée pendant le déjeuner). Seul le \emph{Titre} est
        obligatoire~; auteur, éditeur et notes sont optionnels.
\end{itemize}

\subsection{Parcourir en librairie}

La vue liste se transforme en pile de cartes sur les petits
écrans \string— conçue pour parcourir rapidement d'un pouce dans
une librairie. Chaque carte affiche le titre, l'auteur (s'il est
connu), l'ISBN (s'il y en a un), l'éditeur et la date d'ajout.

Pour imprimer une copie papier ou sauver un PDF pour usage
hors-ligne, cliquez sur \emph{Imprimer} en haut de la liste.
mybibli rend une page d'impression épurée~; utilisez \emph{Imprimer
en PDF} du navigateur (Fichier \string→ Imprimer \string→
Enregistrer en PDF, ou Ctrl/Cmd-P) pour générer un PDF à
emporter à la librairie.

\subsection{Marquer un livre comme acheté}

Quand vous achetez un livre, ouvrez son entrée wish list (ou
tapez \emph{Acheté} directement depuis la ligne de la liste),
confirmez dans la modale, et l'entrée est mise en
soft-delete. L'action est réversible depuis le panneau Corbeille
de l'admin pendant 30 jours.

\subsection{Auto-détection au scan catalogue}

Quand vous cataloguez un livre fraîchement acheté via le flux
standard de scan (\texttt{/catalog}), mybibli vérifie si son ISBN
correspond à une entrée wish list active. Si oui, un bandeau
émeraude apparaît dans la zone de feedback~: \emph{Cet ISBN est
sur votre wish list \string— la retirer~?} Un clic sur le bouton
\emph{Retirer de la wish list} ouvre la même modale \emph{Marquer
comme acheté}, et la confirmation supprime l'entrée. Aucune
suppression silencieuse.

\section{Ce que vaut la bibliothèque}\label{sec:valuation}

Deux montants facultatifs peuvent être enregistrés sur chaque volume :
ce qu'il vous a \emph{coûté} (prix d'achat) et ce qu'il \emph{vaut}
aujourd'hui (valeur actuelle), chacun avec sa devise. Les deux sont
vides par défaut et rien ne force à les remplir --- une bibliothèque qui
n'est pas un placement peut ignorer cette section.

Si vous les renseignez, \texttt{/stats/value}\index{valorisation} en
fait la somme sur tout le catalogue : un total par devise, puis une
ventilation par genre et par série. Les volumes sans valeur sont
ignorés plutôt que comptés à zéro : un catalogue partiellement valorisé
dit donc la vérité sur la partie que vous avez valorisée, au lieu de
diluer la moyenne.

\begin{itemize}
  \item La \textbf{devise par défaut} des nouvelles valorisations se
        règle dans \emph{Administration \textrightarrow{} System} ; ce
        n'est qu'un pré-remplissage, chaque volume peut en porter une
        autre.
  \item La \textbf{carte du tableau de bord} qui affiche le total estimé
        est \emph{désactivée} par défaut. Activez-la au même endroit si
        vous voulez ce chiffre sous les yeux chaque jour --- ce n'est pas
        le souhait de tous les foyers.
\end{itemize}

\begin{note}
L'assurance est la raison habituelle de tenir ces chiffres. Si c'est la
vôtre, notez le prix d'achat autant que la valeur actuelle : le premier
est un fait que vous pouvez prouver, la seconde une estimation qu'on
vous demandera de justifier.
\end{note}
